\documentclass[]{../../../../tex/classes/styledArticle}
\usepackage{../../../../tex/packages/mathShortcuts}
\usepackage{../../../../tex/packages/theoremsSupport}



\newcommand\gl {\mathfrak{gl}}
\newcommand\g  {\mathfrak{g}}
\newcommand\ad {\mathrm{ad}}
\newcommand\Sl {\mathfrak{sl}}
\newcommand\gr {\mathrm{gr}}
\newcommand\End{\mathrm{End}}
\newcommand\Fi {\mathfrak{h}}
\newcommand\Fh {\mathfrak{h}}

\newcommand\asc{\mathscr{A}}

\author{Shahar Perets}
\title{Infinite Dimensional Lie Algebras $\sim$ \textit{Lecture 2}}
\begin{document}
	\maketitle
	
	\textbf{Lectures by Prof. Evgeney Feigin}
	
	
	Last time we discussed about a $\g$-Lie algebra, and imposes a $U(\g)$ the unital associative alg. and we saw the PBW theorem that states that $U(\g) \overset{\gr}{\to}S(\g)$ the symetric representation. In the HW we saw that given $\rho \co \g \to \gl(\g)$ reps. (finite-dim) defines as $\bar \rho \co U(\g) \to \End(V)$ is a reps. of associative algebras, by $x \in \g \to \bar x \in U(\g), U(\g)$ generated by $\bar x$ and then $\bar \rho(\bar x) = \rho(x)$ can be extended to non 1-dim elements using $\bar\rho(x_1 \dots x_j) = \bar \rho(x_1) \cdots \rho(x_j)$. Similarly each $\bar\rho \co U(\g) \to \End(V)$ gives as $\rho \co g \to \gl(V), \rho(x) = \bar\rho(\bar x)$. 
	
	\textbf{Def. \textit{cyclic reps. of $\g$}. }given is a map $\rho \co \g \to \gl(V)$, then we have a reps. $\bar \rho \co U(\g) \to \End(V)$. Then $\rho$ is \textit{cyclic }if $\exists v \in V $ s.t. $U(\g) \cdot v = V = \ccb{\bar\rho(a) \cdot v \mid a \in U(\g)}$. 
	
	\textbf{Ex. }given $g = \Sl_2$ and $\rho_m \co \Sl_2 \to \gl(V_m)$, when $\dim V_m = m + 1$. We saw that $\rho_m(f)$ reduces each vectors after $m + 1$ compositions on $\rho_m(f)$ to $0$. Then for $v$ s.t. it takes the maximal amount of compositions we indeed get that $U(\Sl_2) \cdot v = V = \ccb{\rho_m(f)^{i}v} = \ccb{ \bar \rho_m(f)^{i}v}$ 
%	(intuition: Jordan blocks). 
	
%	[Can't understand the next example]
%	\textbf{Ex. }for $g = \gl_n$ and $\rho \co \gl_n \to \gl()$
	
	\textbf{Remark. }each $\rho \co g \to \gl(v) $ is irreducible reps., then every vector is cyclic, because $U(\g)v \subset V$ is a non-zero subreps. 
	
	There are cyclic non-irreducible reps., e.g.: \textbf{Ex. }($\g = \Sl_2$) $\rho_1 \co \g \to \gl(V_1)$ and $\rho_2 \co \g \to \gl(V_2)$ and then 
	\begin{align*}
		&\overset{v}{\cdot}\to \cdot \quad V_1 \\
		&\overset{v}{\cdot} \to \overset{v}{\cdot} \to {\cdot} \quad V_2
	\end{align*}
	Then: 
	\[ v + u, \quad h(u \oplus u) = v \oplus 2u  \mapsto v, u \in U(\g)(v \oplus u) \]
	Applying $f v, fu , f^{2}u$: everything is cyclic ($(\bar h - 1)(v \oplus u) = u$). 
	
	Fix the std. basis $v_1 \dots v_n$. Consider $\gl_n = n_- + \oplus \Fi \oplus n_+$ where $n_i$ is the lower triangular, $n_+$ the upper tirangular and $\Fi$ is diagonal matrices. Then it has a PBW basis $x_1^{i_1} \dots x_s^{i_s}, y_1^{i_1} \dots y_t^{i_t}, z_1^{e_1} \dots z_r^{\ml_r}$ which is compatible in the above decomposition. 
	
	\textbf{Ex. }$\rho \co \gl_n \to \gl(\Lg^{k}V)$, $v = v_1 \land \dots \land v_k$, then $U(n_+)\cdot v = 0$, because $n_+ = \Sp\{E_{ij} \mid i < j\}$,  and $E_{ij}v_j = v_i$. And $U(\Fi) = \Sp v$ since $h = \diag(x_1 \dots x_n)$ and $\rho(v)v= (x_1 + \dots x_n)v$. Finally $U(n_-)v = \Lg^{k}V$. 
	
	\textbf{Def. }highest weight reps.: $\g = n_+ \oplus \Fi \oplus n_-, \ $ $\Fi$-abelian, $[\Fi, n_-] \subset n_-, [\Fi, n_+] \subset n_+$, then we can define $\rho \co \g \to \gl(W)$ s.t. $\exists w \in W \co \rho(n_+)w = 0, \ \rho(h)w = \lg(h)w \ (h \in \Fi, \lg \in \Fi^*), \ U(n_-)w = W$. 
	
	Take $\gl \supset \Fi$ abelian subalgebra, and suppose we have a reps. $\rho \co \g \to \gl(V)$ s.t. all of the operators $\rho(w)$ commutes. Then there exists a common eigenbasis for all $\rho(h), h \in \Fh$. Let $v_1 \dots v_N \in V$, then $\rho(h)v_i = \lg_i(h)v_i$ for all $\lg_1 \dots \lg_N \in \Fh^*$. In this scenario, we define: 
	
	\textbf{Def. }$\mathrm{ch}_V \cod \sum_{\lg \in \Fh} \dim V(\lg)z^{\lg}$ where: 
	\[ V(\lg) = \ccb{v \in V \co \rho(h)v = \lg(h)v_1 \ \forall h \in \Fh} \]
	
	And $z^{\lg}$ is a symbolic symbol. Let $h_1 \dots h_n$ a basis of $\Fh$. Then we can write $z^{\lg} = z_1^{\lg(h_1)} \cdots z_n^{\lg(h_n)}$. Hence: 
	\[ \mathrm{ch}_V = \sum_{\lg \in \Fh} \dim V(\lg) \underbrace{\prod_{i = 1}^{n} z_i^{\lg(h_i)}}_{ \in \C[z_1 \dots z_n]} \]
	
	\textbf{Ex. }$\g = \gl_n$ and $\Fh = \Sp h$ where $ h = \binom{1 \ \ \ 0}{0 \ -1}$, in respect to $\rho = \rho_m$ (the $m + 1$-dim reps. of $\sl_2$). Then what is $\mathrm{ch}_{V_m}(z_1)$? We got a basis: 
	\[ \overset{v_1}{\cdot} \overset{f}{\to} \overset{v_2}{\cdot} \overset{f}{\to} \cdots \overset{v_{m + 1}}{\cdot} \]
	of $V_{m}$, hence $\rho_m(h)v_i = (m - 2(i - 1))v_i$, and finally: 
	\[ \mathrm{ch}_V = z^{m} + z^{m + 2} + \cdots + z^{-m} \]
	(we have $m + 1$ elements above)
	
	\textbf{Ex. }$\g =  \gl_n$ and $\rho \co \g \to \gl(\Lg^{k}V)$, $\Fh$=diagonal part. Then $h_i = E_{ii}$ ia a basis of $\Fi$, $i >  1\dots n$. What is $\mathrm{ch}_{\Lg^{k}V}(z_1 \dots z_n)$? For $k = 1$ we have $ ch = z_1 + \cdots + z_n$. In the general case: 
	\[ d_{\Lg^{k}V}=(z_1 \dots z_n) = \sum_{1 , i_1 < \cdots < i_k \le n }z_1 \cdots z_n, \] 
	
	Where $v_{i1} \dots\land V_{ik}$ ia a basis of $\Lg^{k}V$. 
	
	\textbf{Remark.} $\mathrm{ch}_{V \oplus W} = \mathrm{ch}_{V } \oplus \mathrm{ch}_{W}$. 
	
	
%	\textbf{Claim. }$\mathrm{ch}_V(...) = \dim V$ 
	
	
	That was a crash course in finite-dim Lie algebras. Now we'll deal with the infinite dimensional case. 
	
	\section*{(Actual) Infinite Dimensional Lie algebras}
	\subsection{The Heisenberg Algebra}
	
	Quite similar to $\Sl_2$, there is an infinite algebra that shows up in many examples, called the \textit{Heisenberg Algebra} or the \textit{Oscillator Algebra}. 
	
	\textbf{Def. }$\asc \cod \Sp\{a_n, n \in \Z\} \oplus \C K$, with the bracket: 
	\[ [K, a_n] = 0 \ \forall n \in \Z \quad\quad [a_m, a_n] = m \dg_{n + m, 0}K \]
	\textbf{Remark}. There exists books where for $n > 0$ uses the notion $a_n = a_n^{+}$ and $a_{-n} = a_n^{-}$. Those are called the Annihilation/Creation operators respectively. 
	
	\textbf{Claim. }$Z(\asc) = \Sp\{a_0, K\}$. 
	
	\textbf{Claim. }$\asc$ is the direct sum of $\asc = \asc_+ \oplus \asc_0 \oplus \asc_{-1}$ where $\asc_0 = Z(\asc) = \Sp\{K, a_0\}$ and $\asc_+ = \Sp\{a_n, n > 0\}, \ \asc_- = \Sp\{a_n, \ n< 0\}$. All of these three subalgebras are commutative. 
	
%	\newcommand\xsc{\mathscr{X}}
	
	\textbf{Def. \textit{Fock Modules}. }Fix $\mu \in \C$. Then $\pi_\mu \co \asc \to \gl(B)$ (where $B = \C[x_1, x_2, x_3, \dots ]$ the algebra of finite-linear-combinations of monomials of the form $x_i^{k_1} \dots x_{i_N}^{x_N}$), is defined as: 
	\[ \pi_\mu(K) = \mathrm{Id} \quad \quad \pi_\mu(a_0) = \mu \mathrm{Id}\quad\quad (n> 0) \ \ \pi_{\mu}(a_n) = \frac{\dd}{\dd x_n}\quad \quad (n < 0) \ \ \pi_\mu(a_n) = -nx_{-n} \]
	
	(sometime for $n > 0$ some authors writes $\pi_{\mu}(a_{-n}) =nx_n$)
	
	Explicitly, let $P \in \C[x_1, x_2 \dots ]$, then for $n > 0$ we establish $\pi_\mu(a_n) P= \frac{\dd}{\dd x_n}P$ and $\pi_\mu(a_n)P = nx_n P$. 
	
	\textbf{Proposition. }$\pi_\mu$ is an irreducible reps. of $\asc$ for any $\mu \in \C$. \begin{proof}
		\begin{itemize}
			First of all we need to prove that this is indeed a reps.
			\item  $\pi_\mu(K)$ and $\pi_\mu(a_0)$ are proportional to $\mathrm{Id}$, so $[\pi_\mu(K), ? ] = [\pi_\mu(a_0), ?] = 0$. 
			\item 		For all of the other cases $[\pi_\mu(a_n), \pi_\mu(a_m)] = 0$ if $m - n \neq 0$ since in such case the lie bracket operates on the diff. variables. 
			\item Assume $m, n$ are the same sign. Then any two multiplications commutes. 
			\item Otherwise: 
			\[ [\pi_\mu(a_m), \pi_\mu(a_{-m})]P = \csb{\frac{\dd}{\dd x_m}, mx_m}P = m\cl{\frac{\dd}{\dd x_m}(x_m P) - x_m\cl{\frac{\dd}{\dd x_m}P}} = mP = \pi_\mu(K)P \]
		\end{itemize}
		It remains to prove that $\pi_\mu$ is irr. We need to show that there are no nontrivial subreps. Let $W \subset B$ a subspace. and let $P \in W$. We need to prove that $W = B$. 
		\begin{itemize}
			\item Let us show that $1 \in W$. Let $x_1^{k_1} \dots x_N^{k_N}$ be a monomial in $P$ of the maximal possible total degree. Then $\pi_\mu(a_1^{k_1}) \cdots \pi_\mu(a_n)^{k_N} P = \cl{\frac{\dd}{\dd x_1}}^{k_1} \cdots \cl{\frac{\dd}{\dd x_N}}^{k_N}P = (\text{const}) 1$ (can be done in the universal enveloping too). Therefore $1 \in W$. 
			\item For any $P \in B$, we know that: 
			\[ P\cl{\pi_{\mu}(a_{-1}), \pi_{\mu}\cl{\frac{a_{-2}}{2}}, \pi_\mu\cl{\frac{a_{-3}}{3}}, \dots} \cdot 1 = P \]
			Therefore $W = B$. 
		\end{itemize}
	\end{proof}

	\textbf{Remark. }$U(\asc) \cdot 1 = B, U(\asc_+) \cdot 1 = 0$ since $\pi(K) 1 = 1 \land \pi(a_n)1 = $. 
	
	\textbf{Proposition. }assume $\pi \co \asc \to \gl(V)$ s.t. there exists $v \in V $ with the following properties: $\asc_+v = 0, \pi(K)v = v, \pi(a_0)v = \mu (v)$. Then: $\pi(a_{-i}^{i_1}, \dots a_{-N}^{i_N})v$ are linearly independent (for different valuees of $a_1 \dots a_n$). If they span $V$, then $\pi = \pi_\mu$. \begin{proof}
		Consider a map $B \to V$ s.t. $1 \mapsto v$ and then extends to homomorphism of reps. If we succeed, the above will follow (because the map will have no kernel). It can be extended via the rule: 
		\[ P \mapsto P\cl{a_1, \frac{a_2}{2}, \frac{a_3}{3} \dots }v \]
		And we should be able to prove that $\pi_\mu(a_n) \cdot P \mapsto \pi(a_n)P(a_1, \frac{a_2}{2} \dots )v$. 
		
		[Finish the proof in HW]
		
		since $\pi_\mu$ is irr., all the monomials $\pi(a_{-1}^{i_1} \dots a_{-N}^{i_N})$ are linearly independent (the kernel is empty). If $B \to V$ is subjective, it is an isomorphism. 
	\end{proof}
	
	Is there any way to define the ''size'' of $B$? We can partition it to parts, namely: 
	\[ \deg\cl{x_1^{k_1}, \dots x_N^{k_N}} = \sum_{i = 1}^{n} ik_i\quad(k \in \Z_{\ge 0})\quad \deg x_i = i \]
	Then we can write: 
	\[ B = \bigoplus_{m \ge 0}B_m \]
	In a composition called \textit{the energy decomposition} (ahhhh physics vibes) with $d \co B \to B$ called \textit{the energy operator} defined as $d(x^{k_1} \dots x_N^{k_N}) = \cl{\sum i k_i} \prod x_i^{k_i}$. Then $B_m \cod \{b \in B \co d(b) = mb\}$. 
	
	How to compute $\sum_{m \ge 0} z^{m}\dim B_m$? 
	
	\textbf{Remark. } some polynomials ay no have concrete degree. But there are some polynomials which are homogeneous. We say that a polynomial is from degree $n$ if each monomial in the sum has degree $n$. Though the energy operator is extendable linearly. 
	
	Back to the computation of $\sum_{m \ge 0} z^{m}\dim B_m$. 
	
	\textbf{Lemma. }$\dim B_m = p(m)$ where $p(m)$ is the number of ways to write $m = \lg_1 + \lg_2 + \cdots + \lg_s$ where $\lg_i \in \Z_{\ge 0}$ and $\lg_1 \ge \lg_2 \ge \cdots \lg_s$. 
	
	E.g. for $m = 4$, $p(4) = 5$ (for $(4, 3 + 1, 2 + 2, 2 + 1 + 1, 1 + 1 + 1 + 1)$). \begin{proof}
		Suppose $\sum^{N}_{i = 1} ik_i = m$, then the decomp: 
		\[ m = \underbrace{N + N + \cdots  + N}_{K_N \ \text{times}} + \underbrace{(N - 1) + \cdots + (N - 1)}_{K_{N - 1}} + \cdots + \underbrace{1 + \cdots + 1}_{K_1} \]
		(Young Diagrams: can't draw). $m$ is the total number of boxes in the Young Diagram. 
	\end{proof}
	
	There is no explicit formula for $p(m)$, but we have some nice properties of the generating series: 
	\[ \sum_{m \ge 0}2^{m}\dim B_m = \sum_{m \ge 0}2^{m}p(m) = \frac{1}{\prod_{i \ge 1}\cl{1 - z^{i}}} = \prod_{i \ge 1}\cl{1+ z^{i} + z^{2i} + z^{3i} + \cdots} \]
	
	\textbf{Remark. }if we want to compute the [..?] of $z^{r}$ them it's enough to calculate $\prod_{i = 1}^{r}\frac{1}{1 - z^{i}}$. \begin{proof}
		Consider the product $\prod_{i \ge 1}(1 + z^{i} + z^{2i} + \cdots)$. Let as [?] the term $z^{i k_i}$ in the  
		[can't understand the handwrite]
	\end{proof}
	
	
	We can define an action of $\asc$ over $B$. $B$ is a module over $\asc$. $B$ is graded via $B_n$-s. Take $a_n \in \asc, b \in B_m$. Then: 
	\[ \pi_\mu(a_n)\cdot b \in B_{-n + m} \]
	\textbf{Lemma. }$\pi_\mu(a_n)$ is of degree $-n$ (as operator on $B$). \begin{proof}
		Assume that $n > 0$. Then $\pi_\mu(a_n) = \frac{\dd}{\dd x_n}$. So
		\[ \pi_\mu(a_m(x_i \dots x_{i_N})) = \begin{cases}
			0 & i \neq m \\
			\text{const} \cdot \frac{x_i \cdots x_{i_N}}{x_n} & \exists a\co [?]
		\end{cases} \]
		\[ \deg\cl{\frac{x_i \cdots x_{i_N}}{x_n}} = \deg\cl{x_{i} \dots x_{i_N}} = -n \]
		
		Similarly, for $n < 0$ $\deg()$
	\end{proof}
	
	
	
	
	%	\ndoc
\end{document}